\subsection{角。}

在本小节及下一小节中，我们假设 A 是一个有序域，从而特征为零。回忆（第六章勘误），若 F 是 A 上一个向量空间，则关系“存在 λ \textgreater{} 0 使 y = λx”是 F - \{0\} 的元素 x、y 之间的一个等价关系，这个关系的每个等价类称为以 0 为原点的开半直线，而开半直线与 \{0\} 的并称为以 0 为原点的闭半直线（或简称半直线）；若 D 是一条直线而 Δ 是含于 D 的一条闭半直线，则 D 是 Δ 与 -Δ 的并，且不含其他闭半直线。我们称一条半直线是迷向的，如果包含它的直线是迷向的。

也回忆（出处同前），给定 A 上有限 n 维的向量空间 F，F 的一个定向就是选取空间 \(\bigwedge^{n}F\) 的两条半直线之一；属于这条半直线的 n-向量称为正的。赋有定向的有限维向量空间称为定向的。

E 的比 \textgreater0 的位似显然构成 H 的一个指数为 2 的子群；我们把它记作 H \(^{+}\) 。典范同态 i : O \(^{+}\) → S \(^{+}\) /H （参见命题 2）是 S \(^{+}\) /H \(^{+}\) 到 S \(^{+}\) /H 上的典范同态与 O \(^{+}\) 到 S \(^{+}\) /H \(^{+}\) 中的典范同态的复合；由于 O \(^{+}\) ∩ H \(^{+}\) = \{1\} ，后一同态是单射。

命题 7. --- 假设 A 是极大有序域且 \(\Phi\) 是正型（§ 7）。则从 O \(^{+}\) 到 S \(^{+}\) /H \(^{+}\) 中的典范同态以及从 O \(^{+}\) /\{1, -1\} 到 S \(^{+}\) /H 中的典范同态都是双射，且 S \(^{+}\) 同构于 O \(^{+}\) × H \(^{+}\) 。

我们已经看到所说的同态是单射，故只需证明第一个是满射。设 \((e_{1}, e_{2})\) 是 E 的一个标准正交基，并设 w 是使 \(w(e_{1}) = e_{2}\) 的正相似变换（第 1 小节命题 1 的推论 1）；于是 \(w^{2} = -1\) （命题 1，b)）。对任一正相似变换 \(u = a + bw\) （ \(a \in \mathbf{A}, b \in \mathbf{A}\) ），我们有 \(\mathrm{N}(u) = a^2 + b^2 > 0\) ，且存在一个且仅一个旋转，它含于 \(\mathbf{A}(\Phi)\) 的与 \(u\) 相同的半直线中，即 \((a^2 + b^2)^{1/2} u\) 。证毕。

推论. --- 给定两条以 0 为原点的半直线 D、D'，存在一个且仅一个旋转 v 使 v(D) = D'。

这些假设确实蕴含 E 不含迷向直线。于是我们的断言由第 1 小节命题 1 的推论 1 得出。

此后我们假设 A 是极大有序域，且型 \(\Phi\) 是正的。在 E 的直线（相应地，以 0 为原点的半直线）的偶 \((\mathrm{D}_1, \mathrm{D}_2)\) 所成的集合中，关系“存在正相似变换（相应地，旋转）u 使 \(u(\mathrm{D}_1) = \mathrm{D}_1'\) 且 \(u(\mathrm{D}_2) = \mathrm{D}_2'\) ”是偶 \((\mathrm{D}_1, \mathrm{D}_2)\) 与 \((\mathrm{D}_1', \mathrm{D}_2')\) 之间的一个等价关系。偶 \((\mathrm{D}_1, \mathrm{D}_2)\) 的等价类按定义称为直线（相应地，半直线）\(\mathrm{D}_1, \mathrm{D}_2\) （依此顺序）的角；我们把它记作 \((\widehat{\mathrm{D}_1, \mathrm{D}_2})\) 。

命题 8. --- 假设 A 是极大有序域，且型 \(\Phi\) 是正的。设 \(D_{1}, D_{2}, D_{1}^{\prime}, D_{2}^{\prime}\) 是 E 中四条以 0 为原点的直线（相应地，半直线）。要使角 \((\widehat{D_{1}, D_{2}})\) 与 \((\widehat{D_{1}^{\prime}, D_{2}^{\prime}})\) 相等，必须且只需角 \((\widehat{D_{1}, D_{1}^{\prime}})\) 与 \((\widehat{D_{2}, D_{2}^{\prime}})\) 相等。

我们来证明所述条件的必要性。设 \(u\) 是使 \(u(\mathrm{D}_{1}) = \mathrm{D}_{1}^{\prime}\) 且 \(u(\mathrm{D}_{2}) = \mathrm{D}_{2}^{\prime}\) 的正相似变换（相应地，旋转）。由命题 1 的推论 3，存在正相似变换（相应地，由命题 7 的推论，存在旋转）\(\nu\) 使 \(\nu(\mathrm{D}_{1}) = \mathrm{D}_{2}\) 。由于群 \(\mathrm{S}^{+}\) （相应地，\(\mathrm{O}^{+}\) ）是交换的，我们有 \(\mathrm{D}_{2}^{\prime} = u(\nu(\mathrm{D}_{1})) = \nu(u(\mathrm{D}_{1})) = \nu(\mathrm{D}_{1}^{\prime})\) ，这表明 \((\widehat{\mathrm{D}_{1}, \mathrm{D}_{1}^{\prime})} = (\widehat{\mathrm{D}_{2}, \mathrm{D}_{2}^{\prime}})\) 。把 \(\mathrm{D}_{2}\) 与 \(\mathrm{D}_{1}^{\prime}\) 互换，充分性即由必要性得出。

由命题 8 得出，对 E 中以 0 为原点的直线（相应地，半直线）的每个角 \((\widehat{\mathrm{D}_1,\mathrm{D}_2})\) ，典范地关联着 \(\mathrm{S}^{+}/\mathrm{H}\) （相应地，\(\mathrm{O}^{+}\) ）中一个完全确定的元素，即正相似变换 \(\nu\) （相应地，旋转 \(\nu\) ）的模 H 的类，这些变换满足 \(u(D_1) = D_2\) （对角的任一代表 ( \(D_1, D_2\) )，角为 ( \(\widehat{D_1, D_2}\) )）。这样我们定义了一个典范双射 \(h\) （相应地，\(h'\) ），它从直线角的集合 \(\mathfrak{A}_0\) （相应地，半直线角的集合 \(\mathfrak{A}\) ）映到 \(S^+/H\) （相应地，\(O^+\) ）上；特别地，对每个 \(\varphi \in \mathfrak{A}\) ，称 \(h(\varphi)\) 为角为 \(\varphi\) 的旋转。对于 \(\mathfrak{A}_0\) 与 \(\mathfrak{A}\) ，我们将通过 \(h^{-1}\) 与 \(h'^{-1}\) 把 \(S^+/H\) 与 \(O^+\) 的交换群结构转移到它们上，并把这样得到的群 \(\mathfrak{A}_0\) 与 \(\mathfrak{A}\) 写成加法。若以 D、\(D', D''\) 表示 E 中以 0 为原点的直线（相应地，半直线），则依定义有

\[
\widehat {(\mathbf {D} , \mathbf {D} ^ {\prime \prime})} = \widehat {(\mathbf {D} , \mathbf {D} ^ {\prime})} + \widehat {(\mathbf {D} ^ {\prime} , \mathbf {D} ^ {\prime \prime})}\tag{9}
\]

（沙勒关系）；

由此推出

\[
\widehat {(\mathrm{D} , \mathrm{D})} = 0, \quad \widehat {(\mathrm{D} , \mathrm{D} ^ {\prime})} = - (\widehat {\mathrm{D} ^ {\prime} , \mathrm{D}}).\tag{10}
\]

注. --- 1) E 中以 0 为原点的直线（相应地，半直线）的集合 L 是交换群 S+/H（相应地，O+）的一个齐性空间，且恒等元是使 L 的所有元素都不变的唯一算子。因此可以对 L 应用第二章 2 \(^{e}\) 版附录 II 第 1 小节的公式；命题 8 是“平行四边形法则”的一个特例，而公式 (9) 与 (10) 是（出处同前）公式 (2) 的特例。

\begin{enumerate}
\def\labelenumi{\arabic{enumi})}
\setcounter{enumi}{1}
\tightlist
\item
  在直线角的群的定义中，可以用典范同构于 S+/H 的群 O+/\{−1, 1\} 来代替群 S+/H（命题 7）。于是从 O+ 到 O+/\{−1, 1\} 上的典范同态对应于从 A 到 A0 上的一个同态，即把两条半直线 Δ、Δ' 的角对应到分别包含 Δ 与 Δ' 的直线 D、D' 的角的那个同态。*当域 A 是实数域时，群 A 是群 A0 的一个二重覆盖。*
\end{enumerate}

由命题 8，凡形如 \((\widehat{\mathbf{D}^{\prime},\mathbf{D}^{\prime\prime}})\) 的直线（相应地，半直线）的角，其中 \(\mathbf{D}'\) 与 \(\mathbf{D}''\) 正交（相应地，凡

形如 \((\widehat{\mathrm{D}}, -\mathrm{D})\) 的角），都相等：这确实由第 1 小节的注 2 得出（相应地，是显然的）。这个直线（相应地，半直线）的角称为直角（相应地，平角）；它是 \(A_{0}\) （相应地，A）中的 2 阶元素。

命题 9. --- 假设域 A 是极大有序的，且 \(\Phi\) 是正型。对每个整数 \(n>1\) ，元素 \(\theta\) （属于直线角的群 \(\mathfrak{A}_{0}\) （相应地，半直线角的群 \(\mathfrak{A}\) ））中满足 \(n\theta=0\) 的个数等于 \(n\) 。

由于 \(\mathfrak{A}_{0}\) 、\(\mathrm{S}^{+}/\mathrm{H}\) 、\(\mathfrak{A}\) 与 \(\mathrm{O}^{+}\) 同构（命题 3），只需对 \(\mathrm{O}^{+}\) 进行证明，即证明恰好有 \(n\) 个旋转 \(\varrho\) 满足 \(\varphi^{n}=1\) 。由于 A 是极大有序域且 \(\mathrm{A}(\Phi)\) 是 A 的次数为 2 的扩张域（命题 1 a)），\(\mathrm{A}(\Phi)\) 是代数闭域（第六章 § 2，第 6 小节，定理 3）。于是 \(n\) 次单位根（在 \(\mathrm{A}(\Phi)\) 中）构成一个 \(n\) 阶循环群（第五章 § 11，第 1 小节，定理 1）。由于 \(\mathrm{N}(u)=u\overline{u}\geqslant0\) （对一切 \(u\in\mathrm{A}(\Phi)\) ），关系 \(u^{n}=1\) 蕴含 \(\mathrm{N}(u)=1\) ，从而 \(u\) 是一个旋转（第 1 小节，注 1）。这就证明了我们的断言。

推论. --- 直角（相应地，平角）是群 \(\mathfrak{A}_{0}\) （相应地，\(\mathfrak{A}\) ）中唯一的 2 阶元素。

最后我们假设平面 E 是定向的。

引理 1. --- 设 u 是 E 的一个正相似变换；所有形如 \(x \wedge u(x)\) 的二向量都属于 \(\bigwedge^{2}\mathbf{E}\) 的同一条闭半直线。

\(u\) 是位似的情形是平凡的。否则，有 \(x \wedge u(x) \neq 0\) （对一切 \(x \neq 0\) ）；设 \(x, y\) 是 \(E\) 的两个向量（ \(x \neq 0, y \neq 0\) ）；存在 \(v \in S^+\) 使 \(y = v(x)\) ，于是 \(y \wedge u(y) = v(x) \wedge uv(x) = v(x) \wedge v(u(x)) = (\det v)(x \wedge u(x))\) ；取 \(E\) 的一个标准正交基，即见 det \(v\) 是正的（命题 1 b)）；由此得我们的断言。

这样，在两个生成元 \(w\) （属于 \(\mathrm{A}(\Phi)\) ，满足 \(w^{2} = -1\) ）之中，存在一个且仅一个，使得二向量 \(x \wedge w(x)\) 对一切 \(x\in\mathbf{E}\) 都是正的。我们选取这个生成元来定义函数 \(c_{w}, s_{w}\) 与 \(t_{w}\) （第 2 小节）。设 \(h\) 与 \(h'\) 是上面定义的从直线角群 \(\mathfrak{A}_{0}\) 到 S+/H 上以及从半直线角群 \(\mathfrak{A}\) 到 O+ 上的典范双射。复合映射 \(t_{w}\circ h\) （从 \(\mathfrak{A}_{0}\) 到射影域 \(\tilde{\mathbf{A}}\) ），\(c_{w}\circ h'\) 与 \(s_{w}\circ h'\) （从 \(\mathfrak{A}\) 到域 A）分别记作 tg、cos 与 sin，称为正切、余弦与正弦函数。映射 \(\varphi\to1/\mathrm{tg}\varphi\) （从 \(\mathfrak{A}_{0}\) 到 \(\tilde{\mathbf{A}}\) ）记作 cotg，称为余切函数。余弦、正弦、正切与余切这些函数统称为三角函数。复合映射 tg∘p 与 cotg∘p（其中 p 表示从 \(\mathfrak{A}\) 到 \(\mathfrak{A}_{0}\) 上的典范同态（上面的注 2））约定仍记作 tg 与 cotg（滥用记号）。

由于这里有 \(\delta = -1\) ，§ 2 的公式 (2)、(8)、(3)、(4)、(5) 与 (7) 给出

\begin{enumerate}
\def\labelenumi{(\arabic{enumi})}
\setcounter{enumi}{10}
\tightlist
\item
\end{enumerate}

\[
\operatorname{tg} \left(\varphi + \varphi^ {\prime}\right) = (\operatorname{tg} \varphi + \operatorname{tg} \varphi^ {\prime}) / (1 - \operatorname{tg} \varphi \operatorname{tg} \varphi^ {\prime})\tag{12}
\]

\[
\operatorname{tg} (2 \varphi) = 2 \operatorname{tg} \varphi / (1 - \operatorname{tg} ^ {2} \varphi)
\]

（对 \(\varphi, \varphi'\) ，\(\mathfrak{A}_0\) 中的）；

\begin{enumerate}
\def\labelenumi{(\arabic{enumi})}
\setcounter{enumi}{12}
\tightlist
\item
\end{enumerate}

\[
\cos^ {2} \theta + \sin^ {2} \theta = 1\tag{14}
\]

\[
\cos (\theta + \theta^ {\prime}) = \cos \theta \cos \theta^ {\prime} - \sin \theta \sin \theta^ {\prime}\tag{15}
\]

\[
\sin (\theta + \theta^ {\prime}) = \sin \theta \cos \theta^ {\prime} + \cos \theta \sin \theta^ {\prime}\tag{16}
\]

\[
\sin (2 \theta) = 2 \operatorname{tg} \theta / (1 + \operatorname{tg} ^ {2} \theta),
\]

\[
\cos (2 \theta) = (1 - \mathrm{tg} ^ {2} \theta) / (1 + \mathrm{tg} ^ {2} \theta)
\]

（对 \(\theta, \theta'\) ，\(\mathfrak{A}\) 中的）。另一方面，依定义或作为前面公式的直接推论，我们有：

\begin{enumerate}
\def\labelenumi{(\arabic{enumi})}
\setcounter{enumi}{16}
\tightlist
\item
\end{enumerate}

\[
\operatorname{tg} \theta = \sin \theta / \cos \theta , \quad \cot \mathrm{g} \theta = \cos \theta / \sin \theta\tag{18}
\]

\[
1 + \mathrm{tg} ^ {2} \theta = 1 / \cos^ {2} \theta , \quad 1 + \cot \mathrm{g} ^ {2} \theta = 1 / \sin^ {2} \theta
\]

（对 θ ∈ 𝒜）。

给定 E 的两个非零向量 \(x, y\) ，这两个向量（依此顺序）之间的角称为并记作 \((\widehat{x}, y)\) ，即它们所属的半直线的角。对 E 的每个向量 \(x\) ，把 A 的一个元素称为 \(x\) 的长度并记作 \(|x|\) ，这个元素就是 \(\Phi(x, x)^{1/2}\) 。

命题 10. --- 假设域 A 是极大有序的，平面 E 是定向的，且型 \(\Phi\) 是正的。对 E 的每对非零向量 \(x\) 、\(y\) ，我们有 \(a\)

\begin{enumerate}
\def\labelenumi{(\arabic{enumi})}
\setcounter{enumi}{18}
\tightlist
\item
\end{enumerate}

\[
\cos \widehat {(x , y)} = \Phi (x, y) / | x |. | y |\tag{20}
\]

\[
\widehat {\sin (x , y)}. e = (x \wedge y) / | x |. | y |,
\]

其中 e 表示使 \(\Phi_{(2)}(e, e) = 1\) 的正二向量。

事实上，由于向量 \(x' = x / |x|\) 与 \(y' = y / |y|\) 的长度都是 1，存在一个且仅一个旋转 \(\nu\) 使 \(\nu(x') = y'\) （第 1 小节，命题 1 的推论 2）。若令 \(\nu = a + bw\) （ \(a, b\) 在 A 中），则依定义有 \(a = \cos(\widehat{x,y})\) 与 \(b = \sin(\widehat{x,y})\) 。关系 \(y' = \nu(x') = ax' + bw(x')\) 给出 \(\Phi(x', y') = a\Phi(x', x') = a\) ，因为 \(x'\) 与 \(w(x')\) 正交（第 1 小节的注 2）；这就证明了 (19)。另一方面，这个关系也给出 \(x' \wedge y' = bx' \wedge w(x') = b.e\) ，这依据开拓 \(\Phi_{(2)}\) （\(\Phi\) 到 \(\bigwedge^{2} E\) 上的）的定义（§ 1，第 9 小节，公式 (37)）以及 \(w\) 的选取；这就证明了 (20)。

注. --- 3) 给定附属於 E 的仿射平面 L 的两条直线（相应地，半直线）D、D'，把 D 与 D' 的角（记作 (D, D')）称为它们的方向在 E 中所成的角（相应地，E 中与 D 及 D' 对应的以 0 为原点的半直线的角）（第二章，2 \(^{e}\) 版，附录 II，第 1 与第 3 小节）。

\begin{enumerate}
\def\labelenumi{\arabic{enumi})}
\setcounter{enumi}{3}
\tightlist
\item
  设 F 是极大有序域 A 上任一维数的一个向量空间，\(\Psi\) 是 F 上一个正的非退化对称双线性型。给定 F 的两个线性无关的向量 \(x\) 、\(y\) ，设 \(\mathbf{F}'\) 是它们张成的向量平面；我们把 \(x\) 与 \(y\) 的角称为把 \(x\) 与 \(y\) 看作平面 \(\mathbf{F}'\) 的元素时的角；记作 \((x,y)\) 。由 (19)，这个角的余弦由下式给出
\end{enumerate}

\[
\cos \widehat {(x , y)} = \Psi (x, y) / | x |. | y |\tag{21}
\]

\((\text{其中 } |x| = \Psi(x, x)^{1/2}\) 也称为向量 \(x)\) 的长度，因而与 \(F'\) 上所选的定向无关；\((x, y)\) 的正弦与正切在改变 \(F'\) 的定向时变号。给定两个成比例的非零向量 \(x, y\) （\(F\) 的），约定 \((x, y) = 0\) 。
% semantic-fragments: legacy-input-seam
\relax{}
